\documentclass[../main.tex]{subfiles}
\begin{document}
    \part{Puissance d'exposant positif}
        \definition{puissance_nombre}
        \begin{exemple}\renewcommand{\arraystretch}{1.5}
            \begin{itemize}[ ]
                \begin{multicols}{2}
                    \item
                    $\begin{array}{ccc}
                        5^3&=&5\times5\times5\\
                        &=&125
                    \end{array}$
                    \item
                    $\begin{array}{ccc}
                        7^2&=&7\times7\\
                        &=&49
                    \end{array}$
                    \columnbreak
                    \item
                    $\begin{array}{ccc}
                        2^4&=&2\times2\times2\times2\\
                        &=&16
                    \end{array}$
                    \item
                    $\begin{array}{ccc}
                        1^3&=&1\times1\times1\\
                        &=&1
                    \end{array}$
                \end{multicols}
            \end{itemize}
        \end{exemple}
        \begin{remarque}
            Par convention $a^0=1$.
        \end{remarque}
        \part{Puissance d'exposant négatif}
            \definition{puissance_nombre_exposant_negatif}
            \begin{exemple}\renewcommand{\arraystretch}{2.2}
                \begin{itemize}[ ]
                    \begin{multicols}{2}
                        \item
                        $\begin{array}{ccc}
                            5^{-3}&=&\dfrac{1}{5^3}\\
                            &=&\dfrac{1}{5\times5\times5}\\
                            &=&\dfrac{1}{125}
                        \end{array}$
                        \item
                        $\begin{array}{ccc}
                            7^{-2}&=&\dfrac{1}{7^2}\\
                            &=&\dfrac{1}{7\times7}\\
                            &=&\dfrac{1}{49}
                        \end{array}$
                        \columnbreak
                        \item
                        $\begin{array}{ccc}
                            2^{-4}&=&\dfrac{1}{2^4}\\
                            &=&\dfrac{1}{2\times2\times2\times2}\\
                            &=&\dfrac{1}{16}
                        \end{array}$
                        \item
                        $\begin{array}{ccc}
                            1^{-3}&=&\dfrac{1}{1^3}\\
                            &=&\dfrac{1}{1\times1\times1}\\
                            &=&1
                        \end{array}$
                    \end{multicols}
                \end{itemize}
            \end{exemple}
    \newpage
    \part{Puissances de 10}
        \subpart{Généralités}
            \definition{puissances_de_dix}
            \begin{exemple}
                \begin{multicols}{2}
                    \item$10^4=10\,000$
                    \item$10^3=1\,000$
                    \columnbreak
                    \item$10^{-5}=0,00001$
                    \item$10^{-2}=0,01$
                \end{multicols}
            \end{exemple}
        \subpart{Notation scientifique}
            \definition{notation_scientifique}
            \begin{exemple}
                \begin{center}
                    \begin{tblr}{   
                        colspec={Q[c,m]Q[c,m]},
                        columns={wd=200pt}
                    }
                        $0,000056=5,6\times10^{-5}$&$3\,904\,560\,000=3,90456\times10^{9}$\\
                        $0,000000903=9,03\times10^{-7}$&$5\,465\,001\,000\,000=5,456001\times10^{12}$
                    \end{tblr}
                \end{center}
            \end{exemple}
        \subpart{Préfixes}
            \begin{center}
                \exported[scale=1.5,lmargin=30,vmargins=30,active=1]%
                <<
                    \begin{tblr}{%
                            colspec={Q[c,m]Q[c,m]Q[c,m]Q[c,m]%
                                Q[c,m]Q[c,m]Q[c,m]Q[c,m]},%
                            hlines,vlines,
                            column{1}={font=\bfseries},
                            row{odd} = {bg=Rainbow1-rouge!40!white},
                            row{even} = {bg=Rainbow1-rouge!70!white}%
                        }
                        Préfixe&\toComplete{giga}&\toComplete{méga}&\toComplete{kilo}&\toComplete{unité}&\toComplete{milli}&\toComplete{micro}&\toComplete{nano}\\
                        Symbole&\toComplete{G}&\toComplete{M}&\toComplete{k}&&\toComplete{m}&\toComplete{\textmu}&\toComplete{n}\\
                        \boldmath$10^\text{n}$&\toComplete{$10^9$}&\toComplete{$10^6$}&\toComplete{$10^3$}&\toComplete{$10^0=1$}&\toComplete{$10^{-3}$}&\toComplete{$10^{-6}$}&\toComplete{$10^{-9}$}
                    \end{tblr}
                >>
            \end{center}
            \begin{exemple}
                \begin{center}
                    \begin{tblr}{   
                        colspec={Q[c,m]Q[c,m]},
                        columns={wd=200pt}
                    }
                        $2,6\times10^{-6}\text{ m}=2,6\text{ \textmu m}$&$9,8\times10^3\text{ m}=9,8\text{ km}$\\
                        $5\times10^{-3}\text{ g}=5\text{ mg}$&$250\times10^9\text{ o}=250\text{ Go}$
                    \end{tblr}
                \end{center}
            \end{exemple}
\end{document}